% =====================================================================
% Appendix G: Internal Color Symmetry as a Modular Subgroup of \Theta
% Status: carrier selection proved, unitary carrier conditional, dynamical gauge field open
% VERSION: v17 Stable Release (updated 2026-02-28: reframed as roadmap)
% =====================================================================

% UBT-AI-PROVENANCE-BEGIN
% schema: ubt-ai-provenance/v1
% tier: B_machine_verified
% ai_assistance: disclosed
% human_review: machine-verification
% editorial_responsibility: Ing. David Jaroš
% policy: ../../AI_PROVENANCE.md
% notice: Machine-verified against named sources or verifiers; individual attestation is not claimed.
% UBT-AI-PROVENANCE-END

\appendix
\section*{Appendix G \\ Internal Color Symmetry as a Modular Subgroup of \texorpdfstring{$\Theta$}{Theta}}
\addcontentsline{toc}{section}{Appendix G: Internal Color Symmetry as a Modular Subgroup of $\Theta$}
\label{app:color_internal_symmetry}

\begin{tcolorbox}[colback=yellow!7!white,colframe=orange!70!black,title=Status: Carrier proved; unitary color dynamics conditional]
The involution construction selects a complex rank-three carrier
$V_c\cong\mathbb{C}^3$ [proved].  Once a unitary structure on that carrier is
postulated, its traceless unitary transformations have Lie algebra
$\mathfrak{su}(3)$ and the standard fundamental/adjoint representation theory
follows.  The selection of the unitary/Yang--Mills dynamics from the canonical
UBT action remains open as \texttt{GAP-SU3-DYN}.

\medskip
\begin{tabular}{@{}p{0.42\linewidth}p{0.19\linewidth}p{0.33\linewidth}@{}}
\hline
Item & Status & Location \\
\hline
Lie algebra $\mathfrak{su}(3)$ on the unitary carrier & \textbf{Conditional} & Theorem~G.A (\S G.16) \\
Fundamental representation on the carrier ($\mathbf{3}$) & \textbf{Conditional} & Proposition~G.B (\S G.16) \\
Adjoint transformation law ($\mathbf{8}$) & \textbf{Conditional} & Theorem~G.C (\S G.16) \\
Confinement & \textbf{Open} & Clay Millennium problem; see \S G.17 \\
Left/right action decoupling & \textbf{Conditional} & Theorem~G.D (\S G.16) \\
\hline
\end{tabular}

\medskip
\textbf{Rank-3 motivation} (\S\ref{sec:G15_rank3}): The fiber rank equals the number of imaginary units of $\mathbb{H}$; the Siegel upper half-space $\mathcal{H}_3$ is the natural parameter space for the internal modular data.
\end{tcolorbox}

\subsection*{G.0 Overview (Core-Compatible, Non-Disruptive)}
This appendix presents a candidate mechanism for QCD color symmetry on a complex rank-three carrier selected by the involution structure of the $\Theta$ field.  It does not identify $\mathrm{SU}(3)$ with the complex-algebra automorphism group of $\mathbb{C}\otimes\mathbb{H}$. The construction preserves UBT core principles:
(i) biquaternionic base for spacetime/kinematics, 
(ii) complex time $\tau=t+i\psi$, 
(iii) the canonical covariant tetrad and its central metric identity,
(iv) gauge/phase data encoded in the holomorphic structure of $\Theta$.
The unitary frame defines a flat color connection. A dynamical connection and its compatibility with the canonical metric and matter sectors remain open.
\textbf{Conditional representation-theory statements are proved in \S G.16 after the unitary-carrier postulate.  Deriving that postulate and the Yang--Mills kinetic term from the canonical UBT action remains open.  Confinement also remains open; see \S G.17.}

\subsection*{G.1 Postulate: Unitary Carrier Dynamics}
Let $\Theta:\; \mathcal{M}\times \mathbb{T}_\psi \to \mathbb{B}\otimes \mathbb{C}$ be the biquaternionic field on spacetime $\mathcal{M}$ with complex time $\tau=t+i\psi$, 
and let $\mathcal{F}$ denote its internal phase manifold. 
\textbf{Postulate (unitary carrier dynamics).}  Having selected the complex carrier $V_c\cong\mathbb{C}^3$ from the involutions, we introduce a unitary frame on that carrier and write
\begin{equation}
\label{eq:G1_theta_factorization}
\Theta(x,\tau)\;=\; \Xi(x,\tau)\,\mathcal{U}(x,\tau),
\end{equation}
where $\Xi$ carries the biquaternionic kinematics and real metric content, while 
$\mathcal{U}(x,\tau)$ is a unitary phase factor acting on a complex rank-$3$ internal fiber:
\begin{equation}
\mathcal{U}(x,\tau)\;\in\;\mathrm{U}(3), 
\qquad \mathcal{U}^\dagger\mathcal{U}=\mathbf{1}_3.
\end{equation}
The \emph{color subgroup} is identified with the traceless part:
\begin{equation}
\label{eq:G1_su3_subgroup}
\mathrm{SU}(3)_{\text{color}}\;\subset\;\mathrm{U}(3), 
\qquad \mathcal{U}=\exp\big(i\,\Phi\big),\quad \Phi^\dagger=\Phi,\quad \mathrm{tr}\,\Phi=0 \;\Rightarrow\; i\Phi\in\mathfrak{su}(3).
\end{equation}
Thus, color rotations are unitary transformations of the selected carrier.  They are module/carrier symmetries and are not, in general, complex-algebra automorphisms of $\mathbb{C}\otimes\mathbb{H}$.


\paragraph{Remark (explicit automorphism counterexample).}
Let the carrier basis be $(\mathbf{I},\mathbf{J},\mathbf{K})$ and consider
\[
U_\theta=\operatorname{diag}(e^{i\theta},e^{-i\theta},1)\in\mathrm{SU}(3).
\]
It sends $\mathbf{I}\mapsto e^{i\theta}\mathbf{I}$,
$\mathbf{J}\mapsto e^{-i\theta}\mathbf{J}$, and
$\mathbf{K}\mapsto\mathbf{K}$.  Although the mixed product
$(e^{i\theta}\mathbf{I})(e^{-i\theta}\mathbf{J})=\mathbf{K}$ survives,
the quadratic relation does not:
\[
(e^{i\theta}\mathbf{I})^2=-e^{2i\theta}\neq-1=\mathbf{I}^2
\]
for generic $\theta$.  Hence this element of $\mathrm{SU}(3)$ is not an
algebra automorphism.  More generally, the complex-algebra automorphism group
of $\mathbb{C}\otimes\mathbb{H}\cong\mathrm{Mat}(2,\mathbb{C})$ is
$\mathrm{PGL}(2,\mathbb{C})$ (inner automorphisms), not $\mathrm{SU}(3)$.
Thus the involutions select $V_c\cong\mathbb{C}^3$ [proved], while the unitary
dynamics on $V_c$ is introduced [open].

\paragraph{Remark (Compatibility).}
Unitarity preserves the Hermitian norm on the introduced carrier.
Compatibility with the canonical metric requires checking its covariant
tetrad readout for the full field representation.

\subsection*{G.2 Modular (Theta-Function) Realization}
A concrete realization uses a multi-variable theta function:
\begin{equation}
\label{eq:G2_multi_theta}
\Theta(x,\tau)\;=\;\sum_{n\in \mathbb{Z}^3}\exp\Big(i\pi\, n^{\!\top}\,\Omega(x,\tau)\,n \;+\; 2\pi i\, n^{\!\top} z(x,\tau)\Big)\, \Xi(x,\tau),
\end{equation}
where $\Omega\in \mathrm{Mat}_{3\times 3}(\mathbb{C})$ is a symmetric period matrix with 
$\mathrm{Im}\,\Omega>0$, and $z\in \mathbb{C}^3$ is the internal phase coordinate. 
Modular transformations $(\Omega,z)\mapsto (\tilde\Omega,\tilde z)$ act on the theta-function data.  Separately, after imposing the unitary-carrier postulate, the traceless part of the local frame is identified with $\mathrm{SU}(3)$ through \eqref{eq:G1_su3_subgroup}.  The implication from the involution-selected carrier to this unitary dynamics is not yet derived from the canonical action.

\paragraph{Interpretation.}
The eight real directions belong to the traceless Hermitian generators
on the unitary carrier. The proposed dictionary from the symmetric period
matrix to all eight directions fails the rank test in \S G.14.

\subsection*{G.3 Flat Color Connection from a Unitary Frame}
For a smooth unitary frame $\mathcal U$ define the Maurer--Cartan form and
its traceless part by
\begin{equation}
\label{eq:G3_MC}
\mathcal A=\mathcal U^{-1}\mathrm d\mathcal U,\qquad
A=\mathcal A-\tfrac13\mathrm{tr}(\mathcal A)\mathbf1_3.
\end{equation}
The identity $\mathrm d(\mathcal U^{-1})=-\mathcal U^{-1}(\mathrm d\mathcal U)
\mathcal U^{-1}$ gives $\mathrm d\mathcal A+\mathcal A\wedge\mathcal A=0$.
Taking the trace gives $\mathrm d\,\mathrm{tr}(\mathcal A)=0$ because the
components of $\mathcal A\wedge\mathcal A$ are commutators. The central
trace term drops out of the commutators, hence
\begin{equation}
\label{eq:G3_curvature}
F=\mathrm dA+A\wedge A=0.
\end{equation}
This frame construction cannot produce generic local Yang--Mills curvature.

\subsection*{G.4 Covariant Derivative on \texorpdfstring{$\Theta$}{Theta} with Color Phase}
Use the right action $\Theta\mapsto\Theta R$ and the connection convention
$A_\mu'=R^{-1}A_\mu R+R^{-1}\partial_\mu R$. The covariant derivative is
\begin{equation}
\label{eq:G4_covD}
D_\mu\Theta=\partial_\mu\Theta-\Theta A_\mu.
\end{equation}
The product rule gives $D_\mu'\Theta'=(D_\mu\Theta)R$.

\paragraph{Yang--Mills comparison.}
On an independently supplied carrier connection the standard kinetic and
matter terms provide a comparison. For real components
$F_{\mu\nu}=-i F_{\mu\nu}^a T^a$ with $\mathrm{tr}(T^aT^b)=\tfrac12\delta^{ab}$,
the usual normalization with coupling $g_s$ gives
\begin{equation}
\label{eq:G4_lagrangian}
\mathcal L_{\mathrm{color}}=
-\frac{1}{4g_s^2}F_{\mu\nu}^aF^{a\mu\nu}
+\mathrm{Re}\langle D_\mu\Theta,D^\mu\Theta\rangle.
\end{equation}
Its curvature term vanishes on the frame ansatz \eqref{eq:G3_MC}.
The origin, normalization and matter embedding of a dynamical connection
are part of \texttt{GAP-SU3-DYN}.

\subsection*{G.5 Algebraic Checks (SU(3) Structure)}
For traceless Hermitian Gell-Mann generators $T^a$ with
$[T^a,T^b]=i f^{abc}T^c$, use $A_\mu=-i A_\mu^aT^a$ with real $A_\mu^a$.
The curvature of a general carrier connection has components
\begin{equation}
F_{\mu\nu}^a=\partial_\mu A_\nu^a-\partial_\nu A_\mu^a
  + f^{abc} A_\mu^b A_\nu^c .
\end{equation}
The component formula is compatible with $F_{\mu\nu}^a=0$ on
\eqref{eq:G3_MC}.

\subsection*{G.5a Explicit Gell-Mann Matrices and Structure Constants}
\label{sec:G5a_gellmann}

For completeness and to enable explicit calculations, we provide the standard $\mathrm{SU}(3)$ generators in the fundamental representation.

\paragraph{Gell-Mann matrices $\lambda_a$.} The eight $3\times 3$ traceless Hermitian matrices generating $\mathfrak{su}(3)$ are ($a=1,\ldots,8$):
\begin{align}
\lambda_1 &= \begin{pmatrix} 0 & 1 & 0 \\ 1 & 0 & 0 \\ 0 & 0 & 0 \end{pmatrix}, \quad
\lambda_2 = \begin{pmatrix} 0 & -i & 0 \\ i & 0 & 0 \\ 0 & 0 & 0 \end{pmatrix}, \quad
\lambda_3 = \begin{pmatrix} 1 & 0 & 0 \\ 0 & -1 & 0 \\ 0 & 0 & 0 \end{pmatrix}, \nonumber \\
\lambda_4 &= \begin{pmatrix} 0 & 0 & 1 \\ 0 & 0 & 0 \\ 1 & 0 & 0 \end{pmatrix}, \quad
\lambda_5 = \begin{pmatrix} 0 & 0 & -i \\ 0 & 0 & 0 \\ i & 0 & 0 \end{pmatrix}, \quad
\lambda_6 = \begin{pmatrix} 0 & 0 & 0 \\ 0 & 0 & 1 \\ 0 & 1 & 0 \end{pmatrix}, \nonumber \\
\lambda_7 &= \begin{pmatrix} 0 & 0 & 0 \\ 0 & 0 & -i \\ 0 & i & 0 \end{pmatrix}, \quad
\lambda_8 = \frac{1}{\sqrt{3}} \begin{pmatrix} 1 & 0 & 0 \\ 0 & 1 & 0 \\ 0 & 0 & -2 \end{pmatrix}.
\label{eq:gellmann-matrices}
\end{align}

The generators are $T^a = \lambda_a/2$, satisfying $\mathrm{tr}(T^a T^b) = \frac{1}{2}\delta^{ab}$.

\paragraph{Commutation relations.} The Lie algebra structure is given by
\begin{equation}
[T^a, T^b] = i f^{abc} T^c,
\label{eq:su3-commutator}
\end{equation}
where $f^{abc}$ are the totally antisymmetric structure constants. Selected non-zero values:
\begin{align}
f^{123} &= 1, \quad f^{147} = f^{165} = f^{246} = f^{257} = f^{345} = f^{367} = \frac{1}{2}, \nonumber \\
f^{458} &= f^{678} = \frac{\sqrt{3}}{2}.
\label{eq:structure-constants}
\end{align}

All other $f^{abc}$ are obtained by antisymmetry: $f^{abc} = -f^{bac} = -f^{acb}$.

\paragraph{Cartan subalgebra.} The maximal commuting subset (Cartan subalgebra) is two-dimensional, generated by $\{T^3, T^8\}$ (diagonal matrices). The eigenvalues of $(T^3, T^8)$ define the weights of representations.

\subsection*{G.5b Explicit Map: Biquaternion Phase Manifold to Cartan Subalgebra}
\label{sec:G5b_phase_map}

\paragraph{Internal phase coordinates.} The internal phase manifold for color is parametrized by three angles $(\phi_1, \phi_2, \phi_3)$ corresponding to the three imaginary quaternion directions in $\Theta$'s internal structure. For the diagonal (Cartan) directions, we focus on the two commuting phases.

\paragraph{Map to Cartan generators.} Define the phase vector:
\begin{equation}
\vec{\phi} = (\phi_1, \phi_2, \phi_3) \in \mathbb{T}^3 \quad \text{(3-torus)}.
\label{eq:phase-vector}
\end{equation}

The traceless diagonal generator in $\mathfrak{su}(3)$ is:
\begin{equation}
\Phi_{\text{Cartan}} = \phi_3 T^3 + \frac{\phi_1 + \phi_2}{\sqrt{3}} T^8,
\label{eq:cartan-map}
\end{equation}
where:
\begin{itemize}
\item $T^3 = \tfrac{1}{2}\mathrm{diag}(1, -1, 0)$ encodes the $\mathrm{U}(1)$ subgroup acting on first two colors
\item $T^8 = \tfrac{1}{2\sqrt{3}}\mathrm{diag}(1, 1, -2)$ encodes the hypercharge-like quantum number
\item The constraint $\mathrm{tr}\,\Phi = 0$ is automatic
\end{itemize}

\paragraph{Full unitary element.} The internal phase frame is:
\begin{equation}
\mathcal{U}(\vec{\phi}) = \exp\left(i\sum_{a=1}^8 \phi^a T^a\right) \in \mathrm{SU}(3),
\label{eq:U-phase-full}
\end{equation}
where $\phi^a$ are eight real phase coordinates on the internal manifold. For small phases, $\mathcal{U} \approx \mathbf{1}_3 + i\sum_a \phi^a T^a$.

\subsection*{G.5c Worked Example: Gauge Potential from Internal Phase and Non-Abelian Isolation}
\label{sec:G5c_worked_example}

\paragraph{Setup.} Consider a specific internal phase configuration:
\begin{equation}
\phi^1(x) = g \, x_1, \quad \phi^2(x) = 0, \quad \phi^3(x) = g \, x_1, \quad \phi^{4-8}(x) = 0,
\label{eq:example-phase}
\end{equation}
where $g$ is a small coupling and $x_1$ is a spacetime coordinate.

\paragraph{Compute color gauge potential.} From Eq.~\eqref{eq:G3_MC}, the Maurer-Cartan form is:
\begin{equation}
A_\mu = \mathcal{U}^\dagger \partial_\mu \mathcal{U} = i\sum_a (\partial_\mu \phi^a) T^a + \mathcal{O}(\phi^2).
\label{eq:A-example-1}
\end{equation}

For the configuration \eqref{eq:example-phase}:
\begin{align}
A_1 &= i g (T^1 + T^3) = \frac{ig}{2} \left[\begin{pmatrix} 0 & 1 & 0 \\ 1 & 0 & 0 \\ 0 & 0 & 0 \end{pmatrix} + \begin{pmatrix} 1 & 0 & 0 \\ 0 & -1 & 0 \\ 0 & 0 & 0 \end{pmatrix}\right] \nonumber \\
&= \frac{ig}{2} \begin{pmatrix} 1 & 1 & 0 \\ 1 & -1 & 0 \\ 0 & 0 & 0 \end{pmatrix}.
\label{eq:A1-example}
\end{align}

All other $A_\mu = 0$ for this configuration.

\paragraph{Compute field strength.} The field strength is:
\begin{equation}
F_{\mu\nu} = \partial_\mu A_\nu - \partial_\nu A_\mu + [A_\mu, A_\nu].
\label{eq:F-example}
\end{equation}

For the configuration \eqref{eq:example-phase}, $A_\mu$ is constant (independent of spacetime coordinates except the linear dependence on $x_1$ in the phase), so all derivatives vanish:
\begin{equation}
\partial_\mu A_\nu = 0 \quad \text{(constant gauge configuration)}.
\label{eq:derivatives-zero}
\end{equation}

Therefore:
\begin{equation}
F_{\mu\nu} = \partial_\mu A_\nu - \partial_\nu A_\mu + [A_\mu, A_\nu] = 0 + 0 + [A_\mu, A_\nu] = 0,
\label{eq:F-vanishes}
\end{equation}
since $A_\mu = A_1 \delta_{\mu 1}$ has only one non-zero component, and $[A_1, A_1] = 0$.

The example illustrates \eqref{eq:G3_curvature}, which holds even for
nonconstant frames. In the standard QCD and QED comparison, electrically
charged quarks can contribute to the one-loop photon polarization with a
color multiplicity of three. Gluons have zero electric charge there. These
comparisons presume the Standard Model charge assignments.

\subsection*{G.6 Embedding in Core UBT and GR Limit}
The carrier comparison preserves its Hermitian norm. A derivation of the
physical metric and its dynamics uses the canonical covariant tetrad and
must check compatibility with the color construction.

\subsection*{G.7 Running, Anomalies, and Consistency (Sketch)}
Standard QCD running and anomaly cancellation are comparisons conditional
on a dynamical connection and the Standard Model matter spectrum.
Commutation of left and right actions is an algebraic result in \S G.16.
Anomalies depend on the chiral spectrum and charge assignments.

\subsection*{G.8 Relation to Multi-Theta (Modular) Data}
The modular parameters $(\Omega,z)$ of \eqref{eq:G2_multi_theta} supply a
candidate parametrization. The proposed period-matrix correspondence is
obstructed by \S G.14. Any parametrization of the smooth unitary frame alone
still gives the flat connection \eqref{eq:G3_MC}.

\subsection*{G.9 Phenomenology and Tests (Program)}
The physical spectrum, strong-coupling normalization and confinement
require a color dynamics beyond this frame construction.

\subsection*{G.10 Speculative Notes (Non-Core, Clearly Marked)}
\emph{Speculative.} If the internal modular space couples weakly to complex-time phase $\psi$, topological defects (domain walls in $(\Omega,z)$) could imprint tiny, energy-dependent modulations in color sector at ultra-high energies. No experimental attempt is known; this is \textbf{not} part of the core claims.

\subsection*{G.11 Summary}
The involutions select a complex rank-three carrier. Its canonical
Hermitian form and complex volume admit an algebraic $\mathrm{SU}(3)$
stabilizer, as shown separately in
\texttt{su3\_stabilizer\_exterior\_fock.tex}. The frame form
\eqref{eq:G3_MC} is flat. The microscopic origin of Yang--Mills dynamics
and matter couplings remains open.

\subsection*{G.12 One-Loop Running in the QCD Comparison}
\label{sec:G12_running}
With the standard independent gauge connection and $n_f$ fundamental Dirac
quark flavors, the QCD comparison has one-loop beta function
\begin{equation}
\label{eq:G12_beta}
\mu\frac{\mathrm d g_s}{\mathrm d\mu}=\beta_0 g_s^3+\mathcal O(g_s^5),
\qquad \beta_0=-\frac{1}{(4\pi)^2}\Big(11-\frac{2n_f}{3}\Big).
\end{equation}
Within one-loop accuracy and for fixed $n_f$ the solution is
\begin{equation}
\alpha_s(\mu)=\frac{\alpha_s(\mu_0)}
{1-8\pi\beta_0\alpha_s(\mu_0)\ln(\mu/\mu_0)},
\qquad \alpha_s=\frac{g_s^2}{4\pi}.
\end{equation}
The normalization of $g_s$ and an induced flow on the modular parameters
are open in UBT.

\subsection*{G.13 Anomaly Checks for the Standard Model Spectrum}
\label{sec:G13_anomalies}
With the Standard Model quark multiplets, the pure color anomaly cancels
between each left-handed $\mathbf 3$ and the corresponding left-handed
conjugate $\bar{\mathbf 3}$. Their cubic traces have opposite signs.

For the mixed $\mathrm{SU}(3)^2\times\mathrm{U}(1)_Y$ anomaly, the quadratic
index is the same for $\mathbf 3$ and $\bar{\mathbf 3}$. Per generation,
the sum of hypercharges over the left-handed colored Weyl multiplets is
\begin{equation}
2Y(Q_L)+Y(u_R^c)+Y(d_R^c)
=2\left(\tfrac16\right)-\tfrac23+\tfrac13=0.
\end{equation}
A single color generator has zero trace, giving a vanishing mixed
gauge--gravitational trace for color. These checks apply with the stated
chiral spectrum and charges. Their realization by $\Theta$ requires the
matter embedding.

\subsection*{G.14 Rank Obstruction for the Period-Matrix Dictionary}
\label{sec:G14_mapping}
The proposed linear variations $\delta\Omega=iS$ with $S$ real, symmetric
and traceless have only five real directions:
\begin{equation}
\dim_{\mathbb R}\operatorname{Sym}_0(3,\mathbb R)=\tfrac{3(3+1)}2-1=5.
\end{equation}
The traceless Hermitian matrices form an eight-dimensional real space.
They decompose as $H=S+iK$, where $S$ is real symmetric and traceless and
$K$ is real antisymmetric. The Gell-Mann directions $T^2,T^5,T^7$ are
antisymmetric under transposition. Taking $iT^a$ in those directions leaves
the space of symmetric period matrices. Even allowing an antisymmetric
addition would give $n^\top K n=0$ in \eqref{eq:G2_multi_theta} for every $n$.

The proposed dictionary $\delta\Omega=i\sum_a\theta^aT^a$ therefore fails
to realize eight independent deformations of the theta function. A
different map involving $z$ requires a separate derivation. Any smooth
unitary frame built by such a map still satisfies \eqref{eq:G3_curvature}.

\subsection*{G.15 Rank-3 Internal Fiber: Motivation from Quaternionic Structure}
\label{sec:G15_rank3}

\paragraph{Three imaginary units of $\mathbb{H}$.}
The biquaternion base algebra $\mathbb{B}=\mathbb{C}\otimes\mathbb{H}$ contains three
independent imaginary units $\mathbf{i},\mathbf{j},\mathbf{k}$ of $\mathbb{H}$ satisfying
$\mathbf{i}^2=\mathbf{j}^2=\mathbf{k}^2=-1$,
$\mathbf{i}\mathbf{j}=-\mathbf{j}\mathbf{i}=\mathbf{k}$, etc.
When $\Theta$ is promoted to carry an internal color fiber, the minimal complex fiber
that realizes all three independent imaginary-unit directions as independent phase degrees
of freedom has rank~3.  A rank-$N$ fiber with $N\neq 3$ cannot simultaneously
(i) carry representations of all three imaginary units independently and
(ii) remain orthogonal to the $\mathrm{U}(1)$ and $\mathrm{SU}(2)$ sectors already
realized on the left.

\paragraph{Siegel upper half-space $\mathcal{H}_3$.}
The multi-variable theta function~\eqref{eq:G2_multi_theta} uses a symmetric $3\times3$
period matrix $\Omega$ valued in the Siegel upper half-space
$\mathcal{H}_3 = \{\Omega\in\mathrm{Mat}_{3\times3}(\mathbb{C}):\Omega^\top=\Omega,\,
\mathrm{Im}\,\Omega>0\}$.
The dimensionality of $\mathcal{H}_3$ is $\dim_\mathbb{R}\mathcal{H}_3 = 3\cdot 4 = 12$
(real and imaginary parts of the 6 independent entries of a $3\times3$ symmetric matrix),
which is the unique genus consistent with three imaginary quaternionic directions:
each diagonal entry $\Omega_{kk}$ ($k=1,2,3$) encodes the phase along $\mathbf{e}_k\in\{\mathbf{i},\mathbf{j},\mathbf{k}\}$,
and the three off-diagonal entries encode mixing.

\paragraph{Phenomenological consistency.}
Rank~3 is also the unique choice consistent with the observed three-color structure of QCD.
Any rank $N\neq 3$ would give $\mathrm{SU}(N)$ color, incompatible with all QCD data.
The algebraic motivation (three imaginary units of $\mathbb{H}$) and the phenomenological
requirement (three quark colors) are thus mutually reinforcing.

% =====================================================================
\subsection*{G.16 Conditional Group, Representation and Covariance Results}
\label{sec:G16_formal_proofs}

\subsubsection*{Theorem G.A (conditional): $\mathfrak{su}(3)$ Lie Algebra on the Traceless Unitary Fiber}

\begin{tcolorbox}[colback=blue!3!white,colframe=blue!40!black,title=Theorem G.A (conditional on the unitary-carrier postulate)]
Let $\mathfrak{h} = \{\Phi\in\mathrm{Mat}_{3\times3}(\mathbb{C}) : \Phi^\dagger=\Phi,\;\mathrm{tr}\,\Phi=0\}$
be the real vector space of $3\times3$ traceless Hermitian matrices, equipped with the bracket
\[[\Phi_1,\Phi_2]_{\mathfrak{su}(3)} \;:=\; -i[\Phi_1,\Phi_2]_{\mathrm{matrix}}.\]
Then $(\mathfrak{h},[\cdot,\cdot]_{\mathfrak{su}(3)})$ is isomorphic to $\mathfrak{su}(3)$.
\end{tcolorbox}

\begin{proof}
\textbf{(1) Dimension.}
A general $3\times3$ Hermitian matrix has $3^2=9$ real degrees of freedom;
the traceless condition $\mathrm{tr}\,\Phi=0$ removes one real parameter, giving
$\dim_\mathbb{R}\mathfrak{h}=8=\dim_\mathbb{R}\mathfrak{su}(3)$.

\medskip
\textbf{(2) Closure under bracket.}
Let $\Phi_1,\Phi_2\in\mathfrak{h}$ and set $C := -i[\Phi_1,\Phi_2]$.
\begin{itemize}
  \item \emph{Hermiticity:}
    $C^\dagger=i[\Phi_2,\Phi_1]=-i[\Phi_1,\Phi_2]=C$.
  \item \emph{Tracelessness:}
    $\mathrm{tr}(C) = -i\,\mathrm{tr}([\Phi_1,\Phi_2]) = -i(\mathrm{tr}(\Phi_1\Phi_2)-\mathrm{tr}(\Phi_2\Phi_1))=0$
    by cyclicity of trace.
\end{itemize}
Therefore $C\in\mathfrak{h}$.

\medskip
\textbf{(3) Jacobi identity.}
$[\Phi_1,[\Phi_2,\Phi_3]_{\mathfrak{su}(3)}]_{\mathfrak{su}(3)} + \mathrm{cyclic}
= (-i)^2[\Phi_1,[\Phi_2,\Phi_3]_{\mathrm{matrix}}]_{\mathrm{matrix}}+\mathrm{cyclic}=0$,
since the matrix commutator satisfies the Jacobi identity.

\medskip
\textbf{(4) Explicit isomorphism.}
The eight Gell-Mann matrices $\lambda_a$ \eqref{eq:gellmann-matrices} are traceless Hermitian
and form a basis of $\mathfrak{h}$.  Setting $T^a=\lambda_a/2$, one verifies by direct computation
(see~\S G.5a) that $[T^a,T^b]=if^{abc}T^c$ with the structure constants $f^{abc}$
of~\eqref{eq:structure-constants}. The map $\Phi\mapsto-i\Phi$ is a real
linear bijection onto the traceless anti-Hermitian matrices and preserves
the bracket $-i[\Phi_1,\Phi_2]_{\mathrm{matrix}}$.
\end{proof}

\begin{tcolorbox}[colback=blue!3!white,colframe=blue!40!black]
\textbf{Corollary.}
$\mathrm{SU}(3)=\{U\in\mathrm{U}(3):\det U=1\}$ has Lie algebra $\mathfrak{su}(3)\cong\mathfrak{h}$.
In UBT notation, $\mathcal{U}=e^{i\Phi}\in\mathrm{SU}(3)$ for any $\Phi\in\mathfrak{h}$.
\end{tcolorbox}

% ---
\subsubsection*{Proposition G.B (conditional): Fundamental Representation $\mathbf{3}$ on the Carrier}

\begin{tcolorbox}[colback=blue!3!white,colframe=blue!40!black,title=Proposition G.B (conditional on the carrier identification)]
The defining column action $v\mapsto\mathcal U v$ of $\mathrm{SU}(3)$ on the carrier
$\mathbb{C}^3 = \mathrm{span}_\mathbb{C}\{|r\rangle,|g\rangle,|b\rangle\}$
is the irreducible fundamental representation $\mathbf{3}$. The identification
of the physical quark sector with this action is a separate assumption.
The Cartan eigenvalues $(T^3,T^8)$ on the three color basis states are:
\[
|r\rangle:\;\bigl(+\tfrac{1}{2},+\tfrac{1}{2\sqrt{3}}\bigr),\quad
|g\rangle:\;\bigl(-\tfrac{1}{2},+\tfrac{1}{2\sqrt{3}}\bigr),\quad
|b\rangle:\;\bigl(0,\,-\tfrac{1}{\sqrt{3}}\bigr).
\]
These are the weights of the fundamental representation $\mathbf{3}$ of $\mathfrak{su}(3)$.
\end{tcolorbox}

\begin{proof}
The defining map $\rho:\mathrm{SU}(3)\to\mathrm{GL}(\mathbb{C}^3)$,
$\rho(\mathcal{U}) = \mathcal{U}$ (standard matrix multiplication on $\mathbb{C}^3$)
is a group homomorphism since $\rho(\mathcal{U}\mathcal{V})=\mathcal{U}\mathcal{V}=\rho(\mathcal{U})\rho(\mathcal{V})$.

\emph{Irreducibility.}
$\mathrm{SU}(3)$ acts transitively on the unit sphere $S^5\subset\mathbb{C}^3$ (for any two unit
vectors $u,v\in\mathbb{C}^3$ there exists $U\in\mathrm{SU}(3)$ with $Uv=u$, by Gram–Schmidt
construction within $\mathrm{SU}(3)$).  Transitivity implies irreducibility~\cite{bröcker2003representations}.
This is the defining representation with Dynkin labels $(1,0)$.
Its conjugate $\bar{\mathbf 3}$ has labels $(0,1)$ and is inequivalent.
Right multiplication on a row $w\mapsto w\mathcal U$ is a right action.
Viewed as a left action on transposed rows by using the inverse, it
gives $w^\top\mapsto\mathcal U^{-T}w^\top$.

\emph{Weight computation.}
$T^3=\frac{1}{2}\mathrm{diag}(1,-1,0)$ acts as $T^3|r\rangle=+\frac{1}{2}|r\rangle$,
$T^3|g\rangle=-\frac{1}{2}|g\rangle$, $T^3|b\rangle=0$.
$T^8=\frac{1}{2\sqrt{3}}\mathrm{diag}(1,1,-2)$ acts as
$T^8|r\rangle=+\frac{1}{2\sqrt{3}}|r\rangle$,
$T^8|g\rangle=+\frac{1}{2\sqrt{3}}|g\rangle$,
$T^8|b\rangle=-\frac{1}{\sqrt{3}}|b\rangle$.
These weight vectors $(\frac{1}{2},\frac{1}{2\sqrt{3}})$, $(-\frac{1}{2},\frac{1}{2\sqrt{3}})$,
$(0,-\frac{1}{\sqrt{3}})$ in the $(T^3,T^8)$ plane are exactly the three weights
of the Dynkin diagram $(1,0)$ fundamental representation of $\mathfrak{su}(3)$.
\end{proof}

% ---
\subsubsection*{Theorem G.C (conditional) Adjoint Covariance on a Unitary Carrier}
\begin{tcolorbox}[colback=blue!3!white,colframe=blue!40!black,title=Theorem G.C (conditional on local SU(3) carrier symmetry)]
Under the right frame change $\Theta\mapsto\Theta R$ with $R$ a smooth
$\mathrm{SU}(3)$ transformation, the connection obeys
\begin{equation}
A_\mu'=R^{-1}A_\mu R+R^{-1}\partial_\mu R.
\label{eq:GC_gauge_transf}
\end{equation}
The curvature and the difference of two connections transform by
conjugation, in the adjoint $\mathbf8$.
\end{tcolorbox}

\begin{proof}
For \eqref{eq:G3_MC}, take the traceless part of
$(\mathcal UR)^{-1}\mathrm d(\mathcal UR)$ using $\det R=1$.
For a general connection the same rule follows from the right derivative
\eqref{eq:G4_covD}. The inhomogeneous terms cancel in a difference.
For the curvature, the identity
$\mathrm d+A'=R^{-1}(\mathrm d+A)R$ on carrier columns gives
$F'=R^{-1}FR$ by squaring. The particular frame connection has $F=0$.
\end{proof}

\subsubsection*{Theorem G.D (conditional) Commuting Left and Right Actions}
\begin{tcolorbox}[colback=blue!3!white,colframe=blue!40!black,title=Theorem G.D (conditional on the stated left/right actions)]
Let $L$ act on the left and $R$ on the right of the stated carrier.
Their actions commute and the derivative
\[
D_\mu\Theta=\partial_\mu\Theta+\Omega_\mu^L\Theta-\Theta A_\mu^R
\]
is covariant under $\Theta'=L\Theta R$.
\end{tcolorbox}

\begin{proof}
Associativity gives $(L\Theta)R=L(\Theta R)$. The map
$(L,R)\mapsto(\Theta\mapsto L\Theta R^{-1})$ represents the direct product.
Its faithful image depends on the kernel, which can include shared
central phases. The global gauge group requires the actual matter actions.

Take $\Omega_\mu^{L'}=L\Omega_\mu^L L^{-1}-(\partial_\mu L)L^{-1}$
and use \eqref{eq:GC_gauge_transf} for $A_\mu^{R'}$. The two product-rule
terms in $\partial_\mu(L\Theta R)$ cancel against the derivatives of $L$
and $R$ in the respective connection transformations. This gives
$D_\mu'\Theta'=L(D_\mu\Theta)R$.
\end{proof}

% =====================================================================
\subsection*{G.17 Confinement Status}
\label{sec:G17_confinement}
Color confinement remains open in UBT. The flat phase connection cannot
provide the nonzero curvature required by the proposed mechanism.
The continuum Yang--Mills existence and mass gap problem is a related
question. The carrier algebra and its adjoint action do not resolve it.

\textbf{Summary.} G.A--G.D establish algebraic and covariant properties
of the introduced carrier. \texttt{GAP-SU3-DYN} still requires a dynamical
connection, its normalization and the matter embedding from the UBT action.

% =====================================================================
% End of Appendix G
% =====================================================================
